\documentclass[11pt]{article}
\usepackage[a4paper,margin=26mm]{geometry}
\usepackage{amsmath,amssymb,amsthm,mathtools,longtable,booktabs,hyperref}
\setlength{\emergencystretch}{2em}
\newtheorem{theorem}{Theorem}
\newtheorem{lemma}[theorem]{Lemma}
\newtheorem{proposition}[theorem]{Proposition}
\newtheorem{remark}[theorem]{Remark}
\newcommand{\Z}{\mathbb Z}
\newcommand{\Q}{\mathbb Q}
\newcommand{\Np}{\mathbb Z_{>0}}
\title{Checkpoint: an algebraic audit of the printed quadratic system:\\57 positive integer unknowns and 45 equations}
\author{\(\mathbb Z\mathbb Q\) Diophantine Project}
\date{23 November 2026}
\begin{document}
\maketitle
\begin{abstract}
We give a complete reversible reduction of the explicit quadratic system of
Mishchenko and Treier from 62 positive integer unknowns and 50 equations to
57 positive integer unknowns and 45 equations. Thus the previously quoted
bound of 57 unknowns and 47 equations follows, with two equations to spare, at the level of the literal printed equations.
The reduction is proved directly; the universality of this particular printed system is not verified here.
A subsequent source comparison found apparent discrepancies in its Pell block;
therefore its transfer bounds remain conditional checkpoints. We record the parameter convention, the full source system,
the reverse map, and a machine-checkable exact algebraic audit. We also
supply a self-contained proof that fifteen integrality conditions and eighteen
positivity conditions suffice for its rational realization.
\end{abstract}

\section{Precisely what is being established}
Fix positive integers \(z,u,y\), which are the source parameters coding a
computably enumerable set, and a positive integer instance \(x\). Put
\[
 a=(2z)^{5^{59}},\qquad \kappa=5^{59}.
\]
Throughout this note, \(x,z,u,y,a,\kappa\) are coefficients, not unknowns.
The degree of an instance polynomial is measured only in its existential
unknowns. This is the convention relevant to undecidability of a class of
polynomial equations. It differs from the convention in an integrality
definition, where the distinguished rational input is a coordinate.

Mishchenko and Treier \cite{MT}, Section~4, present equations (8)--(57) and
state that their positive integer solvability is equivalent to membership
of \(x\) in the set coded by \(z,u,y\). Their source is Jones
\cite{Jones}, equations (D1)--(D37). The construction below does not assume
any version of the \(abc\) conjecture. A transfer to rational undecidability
using a particular definition of \(\Z\) in \(\Q\), however, retains every
hypothesis of that definition.

The frequently quoted pair ``57 unknowns, 47 quadratics'' occurs in the
MathOverflow discussion \cite{MO} as a consequence of Jones's Theorem~3.
That discussion is a lead to the primary source, not the proof used here.
Our argument establishes the stronger pair \((57,45)\) by explicitly
reversible algebraic changes to \cite{MT}.

\begin{remark}[Limit of the source audit]
All fifty rows below were compared with the publicly accessible full text
of \cite{MT}. In particular, the exponent on \(\Gamma_{19}\) in (40) is
\emph{two}. Replacing it by a linear term changes the solution set and is
invalid. The complete 1982 Jones article was not available for a fresh
line-by-line reconstruction of its coding proof during this audit. Thus we
prove the algebraic reduction independently. Universality of this literal
printed system remains the separate unverified hypothesis $\mathsf U_{MT}$. The symbolic checker proves identities and counts,
not universality or the external arithmetic foundations of the project.
\end{remark}

\section{Source discrepancy: promotion is blocked}
The count and affine equivalence below concern the literal printed system.
They do not certify its universality. In particular, (16) and (53) imply
\[
 I^2-(G-1)H^2=1,
\]
whereas Jones's 1980 Theorem~3 has a final Pell equation with coefficient
\(G^2-1\). The square is visible in the original scanned page, independently
of its OCR transcription. The standard Pell-index construction of Matiyasevich
and Robinson (1975), Theorem~4, also has \(G^2-1\). These two equations are
not interchangeable: for \(A=2,E=4,F=7\), the displayed formula
\(G=A+F^2(F^2-A)\) gives \(G=2305\) and \(G-1=48^2\); then
\(I^2-(G-1)H^2=1\) has no integer solution with \(H>0\), by factoring.
The equation with \(G^2-1\) has its usual nontrivial Pell solutions.
This is a counterexample to replacing the local Pell relation, not a
counterexample to the full source universality assertion.

Other discrepancies needing explanation are \(E=iC^2+1\) versus a multiple
of \(C^2\), and the use of \(F^2\) rather than \(D^2\) in the expression
for \(G\). Some changes may be legitimate variants, but they must be
proved, not silently corrected. Accordingly, whenever universality is
invoked below, it is an additional unverified hypothesis \(\mathsf U_{MT}\).
The separate file \texttt{jones1980\_system.py} instead reconstructs the
literal 1980 primary-source theorem and avoids these discrepancies.

\section{The complete source system}
There are 26 symbols \(\Gamma_1,\ldots,\Gamma_{26}\), twenty named symbols
\[
B,D_1,C_1,c,e,\ell,S,T,N,R,P,K,M,U,Y,D,I,E,G,H,
\]
and sixteen further symbols
\[
\alpha,\Delta,\varepsilon,g,t,m,\lambda,\gamma,o,F,i,j,h,w,s,\varphi.
\]
All 62 symbols range over \(\Np\). Upper-case \(G\) and lower-case
\(\gamma\) are distinct from the symbols \(\Gamma_j\).
The following equations reproduce the mathematical system of \cite{MT},
with the source row numbers retained.
\allowdisplaybreaks
\begin{align}
\Gamma_1&=\Gamma_{26}^{2} &&\tag{8}\\
\Gamma_2&=MU &&\tag{9}\\
\Gamma_3&=B(2\Gamma_{23}-B)-1 &&\tag{10}\\
\Gamma_4&=\Gamma_{23}C_1 &&\tag{11}\\
\Gamma_5&=c^2 &&\tag{12}\\
\Gamma_6&=\Gamma_5^2 &&\tag{13}\\
\Gamma_8&=\Gamma_{24}^2 &&\tag{14}\\
\Gamma_9&=\lambda B &&\tag{15}\\
\Gamma_{10}&=GH &&\tag{16}\\
\Gamma_{11}&=F^2 &&\tag{17}\\
\Gamma_{12}&=\Gamma_{23}E &&\tag{18}\\
\Gamma_{13}&=\Gamma_{25}^2 &&\tag{19}\\
\Gamma_{14}&=\Gamma_{23}\Gamma_{25} &&\tag{20}\\
\Gamma_{15}&=N^2 &&\tag{21}\\
\Gamma_{16}&=YK &&\tag{22}\\
\Gamma_{18}&=PK &&\tag{23}\\
\Gamma_{20}&=\Gamma_8\Gamma_{24} &&\tag{24}\\
\Gamma_{21}&=\Gamma_8^2 &&\tag{25}\\
\Gamma_{22}&=\Gamma_6\Gamma_{20} &&\tag{26}\\
B&=2a\Gamma_1^2+1 &&\tag{27}\\
D_1&=\Gamma_{24}+C_1(\Gamma_{23}-B)+\alpha\Gamma_3 &&\tag{28}\\
(\Gamma_4-C_1)(\Gamma_4+C_1)+1&=D_1^2 &&\tag{29}\\
C_1&=\kappa+\Delta(\Gamma_{23}-1) &&\tag{30}\\
c&=1+(\Gamma_{26}-\varepsilon)B+g &&\tag{31}\\
e+2z\Gamma_{26}\ell+2zB\Gamma_6+\Gamma_7&=2z\Gamma_{24} &&\tag{32}\\
\ell&=u+t(B-2z) &&\tag{33}\\
e&=y+m(B-2z) &&\tag{34}\\
S&=g-4z^2\Gamma_{22}+\ell\Gamma_{24}+e(\Gamma_8+4z\Gamma_{22})\notag\\
 &\quad+2z\Gamma_9(-2z\Gamma_{22}+\Gamma_{20}+\Gamma_{21}) &&\tag{35}\\
T&=\Gamma_{24}-1-(\Gamma_{26}-1)\ell
 +(\Gamma_9-2\lambda z)(\Gamma_{24}+\Gamma_8)\notag\\
 &\quad+2z(B-2)\Gamma_{21} &&\tag{36}\\
N&=16z\Gamma_{20}\Gamma_8 &&\tag{37}\\
R&=S(\Gamma_{15}-N)+(T+1)(\Gamma_{15}-1) &&\tag{38}\\
P&=2M\Gamma_2 &&\tag{39}\\
(K-\Gamma_{18})(K+\Gamma_{18})+\Gamma_{19}^{2}&=1 &&\tag{40}\\
(2\Gamma_{25}-2\Gamma_{16}-K)(2\Gamma_{25}-2\Gamma_{16}+K)
 +\Gamma_{17}&=0 &&\tag{41}\\
K&=R+1+h(P-1) &&\tag{42}\\
M&=RY &&\tag{43}\\
U&=\Gamma_{15}w &&\tag{44}\\
Y&=\Gamma_{15}s &&\tag{45}\\
D&=-2\Gamma_{25}-5\gamma+\Gamma_{26}w
       +\Gamma_{23}(\Gamma_{25}+4\gamma) &&\tag{46}\\
I&=D+oF &&\tag{47}\\
(D-\Gamma_{14})(D+\Gamma_{14})+\Gamma_{13}&=1 &&\tag{48}\\
E&=i\Gamma_{13}+1 &&\tag{49}\\
(\Gamma_{12}-E)(\Gamma_{12}+E)-\Gamma_{11}+1&=0 &&\tag{50}\\
G&=\Gamma_{23}+\Gamma_{11}(\Gamma_{11}-\Gamma_{23}) &&\tag{51}\\
H&=2R+1+j\Gamma_{25} &&\tag{52}\\
I^2+H(H-\Gamma_{10})&=1 &&\tag{53}\\
\Gamma_{23}&=\Gamma_2+M &&\tag{54}\\
\Gamma_{24}&=1+\Gamma_9-\lambda &&\tag{55}\\
\Gamma_{25}&=2R+1+C_1+\varphi &&\tag{56}\\
\Gamma_{26}&=\varepsilon+x. &&\tag{57}
\end{align}

\section{A reversible reduction to 57 unknowns}
\begin{lemma}[Affine elimination with a positivity check]
Suppose a polynomial system includes \(v=L(\mathbf t)\), where \(L\)
is affine in the remaining unknowns. Delete this equation and substitute
\(L\) for \(v\) in every remaining equation. If the reduced positive
integer solutions always satisfy \(L(\mathbf t)>0\), restriction and
reconstruction give mutually inverse maps on positive integer solutions.
The degree bound of the equations does not increase.
\end{lemma}
\begin{proof}
Restriction satisfies the substituted rows. Conversely the prescribed
value of \(v\) is an integer, is positive by hypothesis, and satisfies
both the deleted and substituted rows. The two operations are inverse.
Substituting affine polynomials cannot increase total degree.
\end{proof}

\begin{theorem}\label{thm5745}
For each fixed \(x,z,u,y\in\Np\), the source system is equivalent to a
system of 45 quadratic equations in 57 positive integer unknowns. The
transformation is effective in the fixed parameters.
\end{theorem}
\begin{proof}
Make the five simultaneous substitutions
\begin{equation}\label{substitutions}
\begin{aligned}
\Gamma_{23}&=\Gamma_2+M,&
\Gamma_{24}&=1+\Gamma_9-\lambda,\\
\Gamma_{25}&=2R+1+C_1+\varphi,&
\Gamma_{26}&=\varepsilon+x,\\
c&=1+xB+g.
\end{aligned}
\end{equation}
Delete rows (31) and (54)--(57). Row (31) becomes an identity because
\(\Gamma_{26}-\varepsilon=x\). Each right-hand side in
\eqref{substitutions} is affine in the remaining unknowns.
Four reconstructed coordinates are manifestly positive. For the remaining
one, retained rows (15) and (27) give
\[
\Gamma_{24}=1+\lambda(B-1)>0,
\qquad B=2a\Gamma_1^2+1>1.
\]
Thus every positive integer solution of the 45 retained rows reconstructs
all five omitted positive integers and satisfies all fifty source rows.
Conversely, every source solution restricts to a solution of the retained
rows. There are \(62-5=57\) unknowns and \(50-5=45\) equations, each of
degree at most two.
\end{proof}

A complete specification of the reduced equations is: take every source
row except (31), (54), (55), (56), (57), and perform
\eqref{substitutions}. This convention avoids a second fifty-line list
whose only difference is repeated affine expansion. The accompanying
Python module constructs the expanded residuals with the source row
numbers intact.

\begin{remark}[Why the instance parameter must remain a coefficient]
For fixed \(x\), the expression \(1+xB+g\) is affine in \(B,g\).
If \(x\) were itself counted as an unknown, substituting it into (12)
would create the degree-four term \(x^2B^2\). Theorem~\ref{thm5745}
therefore establishes a family of quadratic systems indexed effectively
by \(x\); it does not assert a jointly quadratic presentation in which
\(x\) is also a coordinate. For undecidability this effective family is
exactly what is required.
\end{remark}

Taking the sum of the squares of the 45 residuals gives a polynomial of
degree exactly four in 57 unknowns. It has the same zero set over every
ordered field. The degree is exactly four because at least one residual
has a nonzero degree-two homogeneous part, and a sum of real squares of
those homogeneous parts cannot vanish identically. Under the additional hypothesis \(\mathsf U_{MT}\), the family is undecidable over \(\Np\).
Duplicating two rows gives the weaker bound with exactly 47 rows if that
presentation is desired.

\section{Fifteen integrality conditions suffice}
For this section the 62 source coordinates range over \(\Q\). Impose
integrality only on the following fifteen coordinates:
\begin{equation}\label{fifteen}
\alpha,\Delta,\varepsilon,g,t,m,\lambda,\gamma,o,i,j,h,w,s,\varphi.
\end{equation}
In particular, \(F\) is not in this list.

\begin{lemma}\label{monic}
If \(r\in\Q\) and \(r^d\in\Z\) for an integer \(d\ge1\), then
\(r\in\Z\).
\end{lemma}
\begin{proof}
Write \(r=a/b\) with coprime integers \(a,b\) and \(b>0\).
The condition \(b^d\mid a^d\) forces \(b=1\).
\end{proof}

\begin{proposition}\label{integral}
Every rational solution of all fifty source rows satisfying
\eqref{fifteen} is integral in every coordinate.
\end{proposition}
\begin{proof}
Successively use the source equations to construct integers in the
following order:
\[
\begin{gathered}
\Gamma_{26},\Gamma_1,B,\Gamma_9,c,\ell,e,\Gamma_5,
\Gamma_{24},\Gamma_6,\Gamma_8,\Gamma_{20},\Gamma_{21},
\Gamma_7,\Gamma_{22},T,N,\Gamma_{15},S,R,\\
U,Y,M,\Gamma_2,P,\Gamma_{23},\Gamma_3,C_1,K,
\Gamma_4,\Gamma_{16},\Gamma_{18},D_1,
\Gamma_{25},\Gamma_{13},\Gamma_{14},\Gamma_{17},D,H,E,\Gamma_{12}.
\end{gathered}
\]
Each step is an equality with coefficient \(1\) or \(-1\) on the newly
constructed coordinate and an integer polynomial in coordinates already
known to be integral. Next (50), used in the direction
\[
\Gamma_{11}=(\Gamma_{12}-E)(\Gamma_{12}+E)+1,
\]
gives \(\Gamma_{11}\in\Z\). Equation (17) and
Lemma~\ref{monic} give \(F\in\Z\). Rows (47), (51), (16) then give
\(I,G,\Gamma_{10}\in\Z\). Finally (40) gives
\(\Gamma_{19}^2=1+\Gamma_{18}^2-K^2\in\Z\), and
Lemma~\ref{monic} gives \(\Gamma_{19}\in\Z\).
No division by a variable coefficient has been used.
\end{proof}

\section{Eighteen positivity conditions suffice}
Impose strict positivity on the sixteen primitive coordinates consisting
of \eqref{fifteen} together with \(F\), and also on the two slack
coordinates \(\Gamma_7,\Gamma_{17}\). No separate sign condition is
needed for \(\Gamma_{19}\), because it occurs only squared.

\begin{proposition}\label{positive}
Suppose a rational solution satisfies \eqref{fifteen}, all sixteen
primitive positivity conditions, and
\(\Gamma_7,\Gamma_{17}>0\). After replacing \(\Gamma_{19}\) by its
absolute value if necessary, it is a positive integer source solution.
\end{proposition}
\begin{proof}
Proposition~\ref{integral} first makes every coordinate integral.
Since \(x,\varepsilon\ge1\), \(\Gamma_{26}\ge2\), and hence
\(\Gamma_1\ge4\). Since \(a\ge2z\), (27) gives
\(B\ge64z+1>2z\). Consequently
\[
c=1+xB+g>0,\quad
\ell=u+t(B-2z)>0,\quad e=y+m(B-2z)>z.
\]
Also
\(\Gamma_{24}=1+\lambda(B-1)\ge B\), so all its positive powers
and the powers of \(c\) occurring in (12)--(14), (24)--(26) are positive.

From (32) and \(\Gamma_7>0\),
\begin{equation}\label{slackidentity}
\Gamma_{24}=\Gamma_{26}\ell+B\Gamma_6+
                 \frac{e+\Gamma_7}{2z}>B\Gamma_6.
\end{equation}
Since \(B>2z\),
\[
\Gamma_{20}+\Gamma_{21}-2z\Gamma_{22}
 =\Gamma_{20}(1+\Gamma_{24}-2z\Gamma_6)>0.
\]
Rearrange (35) as
\[
S=g+\ell\Gamma_{24}+e\Gamma_8+4z(e-z)\Gamma_{22}
 +2z\Gamma_9(\Gamma_{20}+\Gamma_{21}-2z\Gamma_{22}),
\]
which is positive. The first part of (36) is positive because
\[
\Gamma_{24}-1-(\Gamma_{26}-1)\ell
 =\ell+B\Gamma_6-1+\frac{e+\Gamma_7}{2z}>0.
\]
The rest of (36) is positive, since
\(\Gamma_9-2\lambda z=\lambda(B-2z)>0\) and \(B>2\).
Thus \(T>0\). Now \(N=16z\Gamma_{24}^5>B\), so \(R>0\) by (38).

Rows (44), (45), (43), (9) give \(U,Y,M,\Gamma_2>0\), and indeed
\(M=RY\ge N^2>B\). Therefore
\[
\Gamma_{23}=\Gamma_2+M>B,
\qquad \Gamma_3=B(2\Gamma_{23}-B)-1>0.
\]
Rows (30), (39), (42) yield \(C_1>0\), \(P\ge2\), \(K>0\).
Then \(\Gamma_4,\Gamma_{16},\Gamma_{18}>0\). Every term on the
right of (28) is positive, so \(D_1>0\); this uses
\(\Gamma_{23}>B\) and does not require a separate Pell sign argument.
Rows (56), (19), (20), (49), (18) give
\(\Gamma_{25},\Gamma_{13},\Gamma_{14},E,\Gamma_{12}>0\).
By (46),
\[
D=(\Gamma_{23}-2)\Gamma_{25}+(4\Gamma_{23}-5)\gamma
      +\Gamma_{26}w>0,
\]
since \(\Gamma_{23}>B\ge65\). Hence \(I=D+oF>0\) and
\(H=2R+1+j\Gamma_{25}>0\). Row (53) gives
\[
H(H-\Gamma_{10})=1-I^2\le0,
\]
so \(\Gamma_{10}\ge H>0\). Equation (16) then gives \(G>0\).
The coordinate \(\Gamma_{11}=F^2\) is positive. Finally,
\[
\Gamma_{19}^2=1+(P^2-1)K^2\ge4,
\]
so its absolute value is a positive integer. Its sign change preserves
every row. This proves positivity of all coordinates.
\end{proof}

\section{A general rational transfer bound}
Assume that a system of \(q\) quadratic equations in \(m\) rational
coordinates, including its distinguished input \(v\), defines exactly
\(v\in\Z\). This is an explicit hypothesis of the following theorem.

For any integer \(n\), Gauss's three-triangular-number theorem gives
\begin{equation}\label{tri}
n>0\quad\Longleftrightarrow\quad
\exists a,b,c\in\Q\quad
 2n-2=a(a+1)+b(b+1)+c(c+1).
\end{equation}
For the converse, each triangular term \(r(r+1)/2\) is at least
\(-1/8\), so \(n\ge5/8>0\). For the forward direction, represent
\(n-1\) as a sum of three triangular numbers with integer arguments.

\begin{theorem}\label{transfer}
Under the stated \((m,q)\) integrality hypothesis and the unverified
source-universality hypothesis \(\mathsf U_{MT}\), there is an effective
family of quartic rational equations with undecidable solvability in at
most
\[
 N=15m+94
\]
unknowns. It is a sum of \(61+15q\) quadratic squares. Every hypothesis
of the chosen integrality definition is retained.
\end{theorem}
\begin{proof}
Start with the 57-coordinate, 45-row reduced system. Attach fifteen copies
of the integrality system at the inputs \eqref{fifteen}, with disjoint
auxiliary witnesses. This costs \(15(m-1)\) coordinates. Attach the
sixteen positivity equations \eqref{tri} for the primitive coordinates,
at a cost of 48 coordinates and sixteen rows.

Each of \(\Gamma_7\) and \(\Gamma_{17}\) appears in exactly one
row, with constant coefficient one. In that row substitute
\[
\Gamma_j=1+\frac12\sum_{k=1}^3 t_{jk}(t_{jk}+1),
\qquad j\in\{7,17\},
\]
then clear the constant denominator. This removes two coordinates and
adds six, while keeping both modified rows quadratic. All reconstructed
source coordinates are integral by Proposition~\ref{integral}; hence
\eqref{tri} and Proposition~\ref{positive} give source soundness.
Completeness follows from integer source witnesses and Gauss's theorem.

The number of coordinates is
\[
(57-2)+48+6+15(m-1)=15m+94,
\]
and the number of residuals is \(45+16+15q\). Sum their squares.
\end{proof}

For the project's 163-coordinate, 137-row integrality interface this
formula gives \(N=2539\), with 2116 quadratic squares. For the
179-coordinate, 151-row interface it gives \(N=2779\), with 2326
squares. For the 56-coordinate, 46-row ordinary-\(abc\) interface it
gives \(N=934\), with 751 squares. These are transfer counts, not an
independent proof of any of the integrality interfaces.

\section{Reproducible algebra and remaining questions}
Run
\begin{verbatim}
python mt_quadratic_system.py --report source_audit_certificate.json
\end{verbatim}
The checker reconstructs all source rows, checks all five substitutions
as polynomial identities, counts variables actually occurring, verifies
the degree bound, tests that the crucial Pell root is squared, and
constructs an acyclic orientation of 45 assignment rows. The five other
rows are (29), (40), (48), (50), (53). The integrality proof uses row (50)
to construct \(\Gamma_{11}\), while the assignment graph instead regards
\(F\) as primitive and uses row (17); these are two different legitimate
orientations and should not be conflated.

No assertion of minimality follows. In particular, a successful reduction
of the number of integrality inputs would save roughly \(m\) coordinates
at once, while further affine substitutions save only a bounded number.
Dropping an integrality condition merely because a rational coordinate
multiplied by an integer is integral is invalid. Likewise, a rational
solution of a Pell conic need not be an integer solution. Any future
compression must preserve the integer reconstruction argument or replace
it with a new one.

\begin{thebibliography}{9}
\bibitem{MT} A. Mishchenko and A. Treier,
\emph{Universal input of knapsack problem for groups},
Journal of Physics: Conference Series \textbf{1050} (2018), 012054.
\href{https://doi.org/10.1088/1742-6596/1050/1/012054}{doi:10.1088/1742-6596/1050/1/012054}.
Mathematical source: Section~4, equations (8)--(57).
Public full text inspected at
\url{https://www.researchgate.net/publication/326636381_Universal_input_of_knapsack_problem_for_groups}.
The article is distributed under CC BY 3.0; this note attributes its source
system and adds independent changes of coordinates and proofs.
\bibitem{Jones} J. P. Jones, \emph{Universal diophantine equation},
Journal of Symbolic Logic \textbf{47} (1982), 549--571.
\href{https://doi.org/10.2307/2273588}{doi:10.2307/2273588}.
\bibitem{Jones80} J. P. Jones, \emph{Undecidable diophantine equations},
Bulletin of the American Mathematical Society \textbf{3} (1980), 859--862.
\href{https://doi.org/10.1090/S0273-0979-1980-14832-6}{doi:10.1090/S0273-0979-1980-14832-6}.
\bibitem{MO} MathOverflow, \emph{Minimum size of undecidable quadratic
Diophantine problems}, question 294432 (2018),
\url{https://mathoverflow.net/questions/294432/minimum-size-of-undecidable-quadratic-diophantine-problems}.
Used only to identify the provenance of the previously quoted 57/47 bound.
\end{thebibliography}
\end{document}
